\documentclass[conference]{IEEEtran}

\usepackage[utf8]{inputenc}
\usepackage{amsmath,amsfonts,amssymb,bm}
\usepackage{graphicx}
\usepackage{booktabs}
\usepackage{multirow}
\usepackage{array}
\usepackage{cite}
\usepackage{xcolor}
\usepackage{tikz}
\usetikzlibrary{arrows.meta,positioning,calc}
\usepackage[hidelinks]{hyperref}

\newcommand{\method}{AudioMimic}
\newcommand{\cof}{Commit Forcing}
\newif\ifshowpapertodos
% Keep evidence gaps visible in the collaborative draft; disable only after review.
\showpapertodostrue
\newcommand{\papertodo}[1]{%
\ifshowpapertodos
\par\smallskip\noindent
{\color{red!75!black}\textbf{TODO:} #1}
\smallskip\par
\fi}

% Research draft updated 2026-09-08; see REVISION_NOTES_20260908.md.
\title{Duet: Discrete–Continuous Motion Generation for Music-Driven Humanoid Dance}

% \author{\IEEEauthorblockN{Tianhu Peng et al.}
% \IEEEauthorblockA{UCL Robotics Institute, University College London}}

\begin{document}
\maketitle

\begin{abstract}
Humanoid dance should be coherent and expressive, appropriate to the music being heard, and realizable by a physical controller. We present \method{}, a robot-native causal dance generation and execution system. Its kinematics-guided discrete--continuous (D+C) representation assigns complementary coarse-structure and fine-detail roles through explicit supervision and jointly decodes both into native robot motion. \cof{} constructs sampled reference transitions and trains continuation from their committed history and reconstructed boundary states. A frozen causal music model supplies timestamped forecasts from arrived audio, and a reference bridge connects committed motion to a fixed SONIC whole-body tracker. Evaluation separates music--motion matching, explicit rhythmic and dynamic adaptation, and dance quality at generation and execution layers. Formal generator comparisons and controlled replay characterize musical matching and downstream tracking limitations. In a two-training-seed codec study, adding the continuous branch reduces structure/detail reconstruction error by approximately 30\%/45\% relative to discrete-only decoding. The proposed supervision raises the discrete-to-continuous structural intervention ratio from 1.20 to 1.41, while continuous-code effective rank reaches 15.1 of 16, supporting a measurable complementary representation.
\end{abstract}

\papertodo{Working draft, not submission-ready. Open evidence gates are indexed
as T01--T12 in TODO.md. Existing pilot results must not be promoted to general
performance, hardware, or safety claims before the corresponding gates close.}
\section{Introduction}
% ================================================================
Humanoid robots are increasingly able to translate external
information into coordinated whole-body behavior.
Language-conditioned motion generation enables robots to perform
actions described in natural language~\cite{jiang2025harmon},
while vision-based control allows them to adapt their movements
to terrain and obstacles~\cite{zhuang2025humanoid}.
These advances motivate richer forms of embodied interaction,
including movement guided by music. Music-responsive dance
requires a robot to transform an evolving auditory signal into
coherent physical expression, with movements that reflect the
soundtrack as it unfolds.

Music provides a distinctive conditioning signal. It typically
does not prescribe a unique action sequence; instead, rhythm,
accents, and dynamic variation guide when and how movements
unfold. Multiple dances can suit the same soundtrack, but
appropriate responses must combine musical correspondence with
coordinated and varied motion. For a physical humanoid, this
relationship must also accommodate the robot's body and remain
perceptible after execution. The objective is therefore to
generate dance whose musical suitability is expressed through
the robot's actual movement.

Music-to-dance models have made substantial progress in
choreographic quality, genre modeling, and whole-body
coordination~\cite{tseng2023edge,li2023finedance,
yang2025megadance,li2025soulnet}.
Recent work has further established important foundations for
robotic performance. DiscoForcing generates motion causally
from streaming audio and transfers human-format motion to a
humanoid through retargeting and tracking~\cite{ji2026discoforcing}.
RoboPerform integrates audio-conditioned expression into a
teacher--student action policy for music-driven dance and
speech-accompanying gestures~\cite{li2026roboperform}.
These approaches demonstrate complementary routes from sound
to embodied behavior.

An open methodological question is how to learn a robot-native
dance reference generator that preserves both whole-body
organization and music-dependent local variation throughout
recursive generation. Generating in a human motion space leaves
the mapping of these properties to the robot to a downstream
retargeting stage. Integrating expression directly into an
action policy couples its realization to the learned controller.
We investigate a complementary design in which musical motion
is planned directly in the robot's native space and delivered
to an independently trained whole-body controller.
This requires more than removing an online conversion step:
the representation must accommodate structural coordination
and expressive detail, and the continuation model must preserve
their consistency as its own outputs become future inputs.
Whether these properties survive tracking is a further
question that must be evaluated at the execution level.

We address the representation problem with a kinematics-guided
discrete--continuous (D+C) design learned from dance data
retargeted to the robot. A discrete vocabulary captures reusable
motion patterns, while a continuous residual supplies variation
beyond quantization. Our supervision explicitly encourages
complementary roles: designated coarse kinematic structure is
supervised through the discrete branch, while the combined
representation reconstructs the full motion. Structural
supervision under residual corruption further encourages this
organization to remain robust to changes in the continuous
component. During generation, both branches receive music
conditions; a discrete planner predicts structure, and a
conditional diffusion model generates the residual given that
structure and the available reference context. Their combination
is decoded into native robot motion by a shared
state-conditioned decoder.

Preserving this organization over time requires training on
the context created by motion commitment. Teacher-forced
histories differ from the model-generated histories encountered
during recursive inference, a challenge addressed broadly by
generated-history and diffusion-forcing
methods~\cite{bengio2015scheduled,chen2024diffusionforcing}.
For our generator, however, a reference transition comprises
discrete codes, continuous residuals, and the boundary state
derived from their decoded motion. We introduce \cof{} to
construct these components jointly during training.
A first pass samples and commits a complete D+C reference
prefix, updating both the latent history and the reference
boundary state. A second pass learns continuation from this
generated context. This procedure exposes the continuation
model to the coupled consequences of its own structural and
detail predictions.

We integrate these components into a causal music-to-robot
system, illustrated in Fig.~\ref{fig:architecture}.
A frozen music forecaster provides predicted conditioning
features from observed audio without accessing future
waveforms. The generator continuously plans and commits
robot-native references, which are delivered to the fixed
SONIC whole-body controller~\cite{luo2026sonic}.
The generator maintains its committed reference context,
while SONIC uses measured robot state for physical feedback
control. Our evaluation connects the proposed mechanisms
to three outcomes: motion quality, musical suitability,
and their retention during execution. Representation
interventions and matched continuation-training comparisons
examine the roles of D+C and \cof{}, while controlled music
pairings and reference--execution comparisons assess whether
the resulting dance responds appropriately to its soundtrack.

Our contributions are:
\begin{itemize}
    \item \textbf{Kinematics-guided robot-native D+C generation:}
    a motion representation and music-conditioned generator
    that establish complementary structural and detail roles
    through explicit kinematic supervision.

    \item \textbf{Commit Forcing for consistent continuation:}
    a training procedure that jointly commits generated
    structure and detail and reconstructs the corresponding
    reference boundary state, enabling continuation learning
    from complete model-generated reference transitions.

    \item \textbf{Causal musical performance with paired
    generation--execution evaluation:}
    a streaming system connecting predicted music conditions
    and robot-native planning to fixed whole-body control,
    together with an evaluation protocol that examines
    musical suitability and motion quality at both layers.
\end{itemize}
\papertodo{T02--T04: close matched learning and music controls before claiming
Commit-Forcing gains, complete D+C superiority, or an injection-architecture winner.}

% ================================================================
\section{Related Work}
% ================================================================

\subsection{Learning Whole-Body Control and Generative Planning}

Reinforcement learning provides an effective framework for
learning humanoid control, while motion imitation supplies
demonstrations of coordinated behaviors that are difficult to
specify through task rewards alone.
PHC develops reusable motion tracking for simulated
humanoids~\cite{luo2023phc}, and ExBody2, BeyondMimic, and
SONIC advance expressive, agile, and scalable control on
humanoid robots~\cite{ji2024exbody2,liao2025beyondmimic,
luo2026sonic}.
These approaches establish the physical execution capabilities
on which higher-level motion generation can build.

Generative models extend this repertoire beyond replaying
individual reference clips.
BeyondMimic combines motion tracking with latent state--action
diffusion, using guidance to compose behaviors and satisfy
downstream objectives~\cite{liao2025beyondmimic}.
An alternative is to retain an explicit motion reference
between planning and control.
Zhang et al.~\cite{zhang2026wholebody} generate terrain-aware
robot references with diffusion, track them using RL, and
fine-tune the tracker with a frozen generator in a closed-loop
setting.
PARC uses physics-based tracking to correct generated motions
and iteratively expand the training dataset for terrain
traversal~\cite{xu2025parc}.
These works demonstrate complementary ways to connect
generative motion models with physical control.
Our work adopts reference-based execution with a fixed
SONIC controller and focuses on music-conditioned robot-native
generation. Commit Forcing updates generated reference
context during training; it does not use physically executed
trajectories to refine the generator or tracker.

\subsection{Motion Representations and Reusable Motor Priors}

Learned motion representations support both generative
modeling and reusable physical skills.
Neural Categorical Priors learn a discrete skill space through
RL-based imitation and use prior shifting to improve
behavioral coverage for downstream control~\cite{zhu2023ncp}.
GPC jointly learns an FSQ-based motion vocabulary and a
control policy through RL, then models skill sequences with
an autoregressive transformer~\cite{shi2026gpc}.
These representations provide latent interfaces to
physics-based control policies.
Our latent representation instead decodes into an explicit
robot motion reference that is executed by a separate tracker.

For kinematic motion generation, T2M-GPT and MoMask use
discrete motion tokens~\cite{zhang2023t2mgpt,guo2024momask},
whereas MLD performs diffusion in a continuous latent
space~\cite{chen2023mld}.
Hybrid approaches connect these modeling choices in different
ways. DisCoRD reconstructs continuous motion from discrete
tokens through rectified-flow decoding~\cite{cho2025discord}.
Most closely, DC-Motion represents human motion using
discrete structural tokens and continuous residual latents,
with masked structural prediction followed by conditional
residual diffusion~\cite{wang2026dcmotion}.
Our work shares the motivation for complementary discrete
and continuous components, while studying their roles in
robot-native streaming generation.
We introduce designated robot-kinematic supervision for the
discrete and residual-corrupted reconstructions, and condition
decoding on the reference boundary state.
The resulting components are generated and committed
jointly, allowing their structural and detail roles to be
examined through controlled interventions.

\subsection{Music Representations, Dance Generation, and Robotic Performance}

Music conditioning ranges from explicit signal descriptors
to learned audio representations.
Features such as onset strength, tempo, and chroma,
available through librosa~\cite{mcfee2015librosa}, provide
interpretable rhythmic and spectral information.
Pretrained representations offer complementary information:
EDGE uses embeddings from the generative music model
Jukebox~\cite{tseng2023edge,dhariwal2020jukebox}, while
Music2Vec, MERT, MusicFM, and MuQ learn transferable
music representations through self-supervised
training~\cite{li2022music2vec,li2023mert,
won2023musicfm,zhu2025muq}.
CLAP, MuLan, and MuQ-MuLan align audio with natural
language, supporting semantic retrieval and
conditioning~\cite{wu2022clap,huang2022mulan,zhu2025muq}.
These approaches provide different sources of acoustic
and semantic information; their suitability for streaming
dance additionally depends on temporal resolution,
accessible audio context, and inference latency.

Music-conditioned dance generation builds on these
representations to model movement and musical correspondence.
FACT and Bailando develop sequence generation and discrete
choreographic modeling~\cite{li2021aist,siyao2022bailando},
while EDGE introduces diffusion-based editable
choreography~\cite{tseng2023edge}.
FineDance expands fine-grained full-body dance
data~\cite{li2023finedance}; LODGE, MEGADance, and SoulNet
address long-form generation, genre-aware modeling,
and coordinated holistic motion,
respectively~\cite{li2024lodge,yang2025megadance,
li2025soulnet}.
These advances establish foundations for expressive dance
generation and music--motion matching.

Generative audio models also offer a route to predictive
conditioning.
MusicGen models compressed audio-token sequences for
controllable music synthesis~\cite{copet2023musicgen},
while Live Music Models and Magenta RealTime~2 support
streaming music generation~\cite{gdmlyria2025live,
magenta2026mrt2}.
In our system, a frozen MRT2 model forecasts music states
from arrived audio, which condition subsequent motion
generation.
This use of predicted states makes audio visibility and
the remaining forecast lead explicit parts of the
music-to-motion interface.

Conditioned motion generation has further reached
physical humanoids.
HARMON generates robot whole-body motion from language
descriptions~\cite{jiang2025harmon}.
RoboPerform uses audio-conditioned style within a
teacher--student control framework for dance and
speech-accompanying gestures~\cite{li2026roboperform}.
DiscoForcing connects causal streaming audio to
human-format motion generation and humanoid execution
through retargeting and tracking~\cite{ji2026discoforcing}.
Our approach performs retargeting during data preparation
and learns a robot-native D+C reference generator, with
music conditioning both discrete planning and continuous
residual generation.
We study musical suitability and motion quality through
paired generation--execution evaluation with a fixed
whole-body controller.

\subsection{Streaming Generation and Training with Generated Context}

Recursive generation introduces a mismatch between
dataset-provided training histories and model-generated
inference histories.
Scheduled Sampling progressively substitutes model predictions
for observed sequence inputs~\cite{bengio2015scheduled}.
Professor Forcing uses adversarial training to align recurrent
dynamics between teacher-forced and free-running
modes~\cite{lamb2016professor}.
Diffusion Forcing trains sequence models with independent
per-token noise levels, supporting causal generation and
flexible conditioning on partially noised
histories~\cite{chen2024diffusionforcing}.
In motion synthesis, MotionStreamer combines causal motion
compression with autoregressive diffusion and a Two-Forward
strategy that replaces a subset of ground-truth latents with
first-pass predictions~\cite{xiao2025motionstreamer}.
DiscoForcing further develops temporal scheduling and
history guidance to balance continuity with responsiveness
to changing audio~\cite{ji2026discoforcing}.

Commit Forcing addresses a specific reference-construction
problem within this broader setting.
A committed D+C transition contains discrete structure,
continuous detail, and a boundary state reconstructed from
their jointly decoded motion.
We sample the structural and residual components with the
deployment samplers, commit their shared prefix, and train
the next continuation using the updated reference context.
This explicitly links history replacement to motion decoding
and boundary-state reconstruction.
The feedback remains within the reference-generation process;
measured physical state is used by the downstream controller.
% ================================================================
\section{Methodology}
\label{sec:methodology}

\subsection{Overview and Problem Formulation}

Our framework generates robot-native dance references
from causally available music and executes them through
a fixed whole-body controller.
Figure~\ref{fig:architecture} summarizes representation
learning and online generation and execution.
Training proceeds in two stages: we first learn a
kinematics-guided discrete--continuous (D+C) motion
codec from retargeted human dance, and then train a
music-conditioned generator with Commit Forcing to
learn continuation from generated reference contexts.

At planning cycle $k$, the generator receives committed
latent history $\mathcal{H}_k$, reference boundary state
$b_k$, and predicted music condition $m_k$.
Let $P_k=\{(q_{0,k,i},r_{k,i})\}_{i=1}^{H}$ denote
an $H$-step latent plan, where $q_0$ is a discrete
code index and $r$ is a continuous residual.
We sample
\begin{equation}
    P_k \sim
    p_\theta(\cdot\mid\mathcal{H}_k,b_k,m_k).
    \label{eq:planning_objective}
\end{equation}
Only the first $C<H$ steps are committed.
Their decoded motion and latent codes update the
reference context before replanning.
The boundary state is derived from committed references;
physical state feedback is handled by the downstream
controller.

\begin{figure*}[t]
    \centering
    \includegraphics[width=\textwidth]
    {figures/mainline_architecture.png}
    \caption{
    \textbf{Overview of the proposed framework.}
    \textbf{(A)} FineDance human motion is retargeted to
    Unitree G1 to learn a kinematics-guided D+C
    representation. Discrete embeddings and continuous
    residuals share a state-conditioned decoder.
    \textbf{(B)} Predicted music states from frozen MRT2
    condition discrete planning and residual diffusion
    through FrameFiLM. The generator, trained with Commit
    Forcing, commits a reference prefix and updates its
    reference context. A bridge delivers references to
    frozen SONIC, whose physical feedback loop uses
    measured robot state.
    }
    \label{fig:architecture}
\end{figure*}

\subsection{Kinematics-Guided Discrete--Continuous
Motion Representation}
\label{sec:dc_representation}

\paragraph{Robot-native motion.}
We use paired music and human dance from
FineDance~\cite{li2023finedance}, with offline motion
retargeting to obtain Unitree G1 trajectories.
Each motion frame is represented by
\begin{equation}
    x_n =
    [\Delta p^{\mathrm{local}}_{xy},\,h,\,
    \sin\Delta\psi,\,\cos\Delta\psi,\,
    \boldsymbol{\vartheta}],
    \label{eq:native_motion}
\end{equation}
where $\Delta p^{\mathrm{local}}_{xy}$ is heading-local
planar displacement, $h$ is root height, $\Delta\psi$
is yaw change, and
$\boldsymbol{\vartheta}\in\mathbb{R}^{29}$ contains
joint angles.
This 34-dimensional representation describes motion
relative to the evolving root frame.
Global planar position and heading are recovered
through deterministic $\mathrm{SE}(2)$ accumulation.

To condition continuation on the preceding reference,
we define the boundary state
\begin{equation}
    b =
    [h,\,v^{\mathrm{local}}_{xy},\,v_z,\,
    \dot{\psi},\,\boldsymbol{\vartheta},\,
    \dot{\boldsymbol{\vartheta}}].
    \label{eq:boundary_state}
\end{equation}
It contains pose and first-order motion information
and conditions both encoding and decoding.

\paragraph{Discrete--continuous decomposition.}
A state-conditioned encoder maps a motion segment
to a latent sequence $z=E_\phi(x,b)$ with a temporal
downsampling factor of two.
Each latent vector is decomposed using a learned
codebook:
\begin{equation}
\begin{split}
    q_{0,i}
    &=\arg\min_j\|z_i-e(j)\|_2^2,\\
    r_i
    &=z_i-\operatorname{sg}(e(q_{0,i})),
\end{split}
\label{eq:dc_decomposition}
\end{equation}
where $e(j)$ is a codebook embedding and
$\operatorname{sg}$ denotes stop-gradient.
The discrete embeddings and residuals are combined
and decoded jointly:
\begin{equation}
    \hat{x}=D_\psi(e(q_0)+r,b).
    \label{eq:shared_decoding}
\end{equation}
The resulting representation couples a finite motion
vocabulary with continuous reconstruction capacity.

\paragraph{Kinematics-guided supervision.}
Full-motion reconstruction alone does not assign
explicit kinematic responsibilities to the two
components.
We therefore supervise discrete reconstruction on
a designated structural view while retaining
full-motion supervision for their joint reconstruction.

The full view $\Phi_f$ contains native motion,
boundary-aware velocities, current-root-relative
link and sole positions, and their velocities.
The structural view $\Phi_s$ selects root, leg,
and waist channels together with torso, wrist,
and sole geometry.
It specifies which aspects of whole-body organization
the discrete component is encouraged to preserve.

Let $e=e(q_0)$ and $r=z-\operatorname{sg}(e)$.
We construct three training paths:
\begin{align}
    \hat{x}_{f}
    &=D_\psi(\operatorname{sg}(e)+r,b),
    \label{eq:full_path}\\
    \hat{x}_{s}
    &=D_\psi(z+\operatorname{sg}(e-z),b),
    \label{eq:structure_path}\\
    \hat{x}_{r}
    &=D_\psi(\operatorname{sg}(e)+\mathcal{C}(r),b),
    \label{eq:robust_path}
\end{align}
where $\mathcal{C}$ removes the residual or adds
small channel-scaled noise.
The discrete-only path uses a straight-through
estimator, while the corrupted-residual path
encourages structural reconstruction to remain
robust to residual variation.

For view $v\in\{f,s\}$, define
\begin{equation}
    d_v(\hat{x},x)
    =
    \frac{1}{|\mathcal{B}_v|}
    \sum_{a\in\mathcal{B}_v}
    \operatorname{MSE}
    \bigl(
    \bar{\Phi}_{v,a}(\hat{x}),
    \bar{\Phi}_{v,a}(x)
    \bigr),
    \label{eq:kinematic_distance}
\end{equation}
where $\mathcal{B}_v$ indexes the native-motion,
velocity, kinematic-position, and kinematic-velocity
blocks of that view.
The bar denotes normalization using training-set
statistics.
Averaging within each block before combining them
prevents high-dimensional geometry from dominating
the objective.

The codec loss is
\begin{equation}
\begin{split}
    \mathcal{L}_{\mathrm{codec}}
    ={}&d_f(\hat{x}_f,x)
    +\lambda_s d_s(\hat{x}_s,x)\\
    &+\lambda_r d_s(\hat{x}_r,x)
    +\mathcal{L}_{\mathrm{VQ}},
\end{split}
\label{eq:codec_objective}
\end{equation}
where $\mathcal{L}_{\mathrm{VQ}}$ comprises codebook
and commitment losses~\cite{vandenOord2017vqvae}.
This objective encourages discrete codes to preserve
designated structure and residuals to provide
complementary detail, without imposing independent
decoding of the two components.

\subsection{Causal Music-Conditioned Motion Generation}
\label{sec:music_generation}

\paragraph{Predicted music states.}
Arrived audio is tokenized with
SpectroStream~\cite{li2025spectrostream} and supplied
to a frozen Magenta RealTime~2
model~\cite{magenta2026mrt2}.
We roll the observed prefix forward and retain
14 predicted temporal music states, denoted FMS14.
These conditions are computed without accessing
future waveform samples.

\paragraph{Timing-aware conditioning.}
The forecaster and generator operate asynchronously.
Each predicted state retains its native timestamp,
so the nominal prediction horizon can be distinguished
from the lead remaining when a reference is due.
For forecast-state timestamp $\tau_j$ and reference
delivery deadline $d_k$, expressed on the same
audio-time axis, the effective lead is
\begin{equation}
    \ell_{k,j}=\tau_j-d_k.
    \label{eq:effective_lead}
\end{equation}
The generator selects the newest completed forecast
and masks states whose effective lead has expired.
The resulting music condition $m_k$ is supplied
to both motion branches.

We use FrameFiLM to introduce this condition through
feature-wise modulation~\cite{perez2018film}.
For an intermediate motion feature $u$, the
modulation has the form
\begin{equation}
    \operatorname{FiLM}(u;m_k)
    =
    \gamma(m_k)\odot u+\beta(m_k),
    \label{eq:music_film}
\end{equation}
where $\gamma$ and $\beta$ are learned conditioning
mappings.
Music thus influences both discrete planning and
continuous residual generation.

\paragraph{Discrete planning and residual diffusion.}
Let $c_k=(\mathcal{H}_k,b_k,m_k)$ denote the shared
generation context.
We factorize the plan distribution as
\begin{equation}
\begin{split}
    p_\theta(q_{0,1:H},r_{1:H}\mid c_k)
    ={}&
    \prod_{i=1}^{H}
    p_{\theta_D}(q_{0,i}\mid q_{0,<i},c_k)\\
    &\cdot
    p_{\theta_C}(r_{1:H}\mid q_{0,1:H},c_k).
\end{split}
\label{eq:generator_factorization}
\end{equation}
A causal Transformer samples the discrete sequence
autoregressively.
A conditional diffusion model then generates the
residual sequence using the sampled structure
and the shared context.
The discrete branch uses cross-entropy supervision.
The residual branch predicts clean residuals from
noised inputs and uses DDIM
sampling~\cite{song2021ddim} at inference.
The sampled embeddings and residuals are combined
through Eq.~\eqref{eq:shared_decoding} to produce
one robot-native reference plan.

\subsection{Commit Forcing for Streaming Generation}
\label{sec:commit_forcing}

\paragraph{Generated reference transitions.}
Streaming inference continues from previously
sampled structure, residuals, and their decoded
boundary state.
Teacher-forced training instead provides dataset
reference contexts.
Commit Forcing exposes the generator to a complete
generated reference transition before supervising
its continuation.

\paragraph{First pass: sample and commit.}
A stop-gradient first pass samples an $H$-step
D+C plan with the deployment sampling procedures.
Its first $C$ discrete codes and corresponding
residuals form one committed prefix.
The reference context is updated as
\begin{equation}
    (\tilde{\mathcal{H}}_{k+1},\tilde b_{k+1})
    =
    \mathcal{U}
    (\mathcal{H}_k,b_k,\tilde P_k^{1:C}),
    \label{eq:commit_update}
\end{equation}
where $\mathcal{U}$ appends the committed latent
prefix, retains the prescribed history length,
and reconstructs the next boundary state from
the decoded reference.
The unused plan suffix is discarded.

A training curriculum selects either a complete
generated transition or a complete dataset
transition, with the probability of selecting
the generated transition increasing to one.
Discrete codes, residuals, and their corresponding
boundary state are selected together.

\paragraph{Second pass: learn continuation.}
The second pass uses the updated reference context
and the music condition for the next horizon.
The discrete branch is supervised against the
next dataset code sequence.
A structural sequence is also sampled
autoregressively from the updated context to
condition residual training.

The dataset residual is defined relative to its
dataset code embedding.
When conditioning on a sampled code, we therefore
rebase the residual target.
Let $z_i^*=e(q_{0,i}^*)+r_i^*$ be the next dataset
target latent and $\tilde q_{0,i}$ the sampled
continuation code.
The rebased target is
\begin{equation}
    \tilde r_i^*=z_i^*-e(\tilde q_{0,i}),
    \label{eq:residual_rebasing}
\end{equation}
so that the sampled embedding and its residual
target sum to the complete target latent.
After residual normalization, the continuation
objective combines classification and denoising:
\begin{equation}
    \mathcal{L}_{\mathrm{gen}}
    =
    \mathcal{L}_{\mathrm{CE}}
    +\lambda_{\mathrm{diff}}\mathcal{L}_{\mathrm{diff}}.
    \label{eq:generator_objective}
\end{equation}

Each training window contains one sampled reference
transition followed by a supervised continuation.
The continuation retains the dataset target latent
under the updated boundary state; no
boundary-dependent re-encoding is performed.
Consequently, rebasing preserves the target latent,
while continuity across generated boundaries is
learned through state conditioning and exposure
to generated reference contexts.

\subsection{Robot Reference Integration and Execution}
\label{sec:robot_execution}

The generator operates in robot-native reference
space, and a dedicated bridge converts committed
motion into inputs for the pretrained SONIC
whole-body controller~\cite{luo2026sonic}.
The bridge accumulates heading-local root
displacements and yaw changes, maps joint references
to the controller convention, and resamples the
30-Hz motion onto a continuous 50-Hz reference grid.
Only committed motion is delivered; the uncommitted
plan suffix is excluded.

SONIC remains fixed and combines the consumed
reference with measured robot state to produce
control actions.
The generator separately updates its latent
history and boundary state from committed reference
motion.
Thus, reference continuation and physical tracking
use distinct feedback paths: generated references
condition subsequent planning, while measured
execution closes the controller loop.
The reference delivery and controller configuration
are held fixed in paired execution comparisons.

% ================================================================
\section{Discussion}
% ================================================================
\textbf{What is supported.} The representation ablation supports complementary
structure/detail roles. The matched replay demonstrates sensitivity to
consumption age even when condition content is held fixed. Tracking the same
recorded trajectories reveals another loss of motion dynamics. Improving
condition timing therefore does not, by itself, establish preserved physical
expression; conversely, ideal-age cached conditions do not eliminate tracking
attenuation.

\textbf{What remains unresolved.} The age-attribution study uses one song and one
sampling seed with an artificial empty-history start. Its trajectory RMSE is
not an aesthetic score. Physical tracking results are MuJoCo measurements,
not a real-hardware validation or safety guarantee. The learned standing-start
candidate has not passed its generation/execution gates. FineDance formal
comparisons, smaller mechanism pilots, and the two-training-seed codec study
have different cohorts and cannot be pooled. Additional architecture and
Commit-Forcing ablations, a calibrated cross-dataset evaluator, and human
judgments are still needed.

\textbf{Research direction.} We prioritize startup train--deployment consistency
and deadline-effective music conditioning before adding model branches. If
tracking-aware reference adaptation is pursued, it should preserve validated
musical events under a fixed controller and be compared with simple smoothing
and clipping. The controller remains fixed in the present study. The final held-out set must be
separate from songs repeatedly used to tune the present diagnostics.

\papertodo{T11/T12: obtain real-G1 records only if hardware claims are retained;
separate edited showcases from continuous trials and report interventions. Finalize
authorship, primary-source citations, figures, venue format and page budget.
Disable visible TODOs only after each retained claim has evidence or is narrowed.}

% ================================================================
\section{Conclusion}
% ================================================================
We presented \method{} toward coherent, expressive robot dance that responds
appropriately to music as it arrives. Kinematics-guided D+C modeling assigns
complementary coarse-structure and fine-detail roles in native robot motion;
\cof{} trains continuation from jointly committed generated references and
their boundary states. Timestamped causal forecasts and a reference bridge
connect this generator to fixed SONIC execution. The evaluation separates MMR,
explicit music adaptation, and dance quality/aesthetics, with runtime and
tracking diagnostics identifying where intended behavior is lost. Existing
representation and layered studies support specific mechanisms, not a claim
that online executed dance already matches offline animation. Matched learning
ablations, learned-startup validation, and music-fit and aesthetic judgments at
both generation and execution remain necessary for that broader conclusion.

\appendices
\section{Detailed Evaluation Definitions and Coverage}
\label{app:evaluation}
\begin{table*}[!t]
\centering
\caption{Three complementary evaluation blocks. R: reported in existing formal
or diagnostic cohorts; I: implemented or exploratory, but final validation or
coverage is incomplete; P: planned. R does not mean that every model, song, and
execution layer has been evaluated. No aggregate quality score is used.}
\label{tab:evaluation_matrix}
\small
\renewcommand{\arraystretch}{1.16}
\setlength{\tabcolsep}{0pt}
\begin{tabular}{@{}>{\raggedright\arraybackslash}p{0.15\textwidth}
                    @{\hspace{0.02\textwidth}}
                    >{\raggedright\arraybackslash}p{0.36\textwidth}
                    @{\hspace{0.02\textwidth}}
                    >{\raggedright\arraybackslash}p{0.45\textwidth}@{}}
\toprule
\textbf{Block} & \textbf{Metrics and direction} & \textbf{Question and evidence status} \\
\midrule
\textbf{1. MMR} & MMR-MS, MMDist $\downarrow$; audio-to-motion and motion-to-audio R@K $\uparrow$, median/mean rank $\downarrow$ &
Does an independent learned evaluator associate the motion with its music?
\textbf{R:} FineDance matching/retrieval subsets. \textbf{I:} cross-dataset calibration;
not a direct aesthetic judgment. \\
\midrule
\textbf{2. Music adaptation} & BAS $\uparrow$; beat-event coverage and timing; onset--motion correlation and response lag &
Does the movement respond to the soundtrack's rhythm and dynamics?
\textbf{R:} BAS and beat diagnostics in formal cohorts; onset/impact diagnostics in
layer analysis. Event coverage is not a mandatory one-to-one beat assignment. \\
& Tempo/phase, section response, style/affect fit; correct versus swapped-audio preference &
\textbf{I:} tempo/style diagnostics. \textbf{P:} validated phrase/affect endpoints
and blinded cross-paired music-fit study at generation and execution. \\
\midrule
\textbf{3. Dance quality and aesthetics} & FID$_k$/FID$_g$/FID$_L$ $\downarrow$;
Div$_k$/Div$_g$/Div$_L$ and same-song multimodality $\leftrightarrow$GT &
Is the motion distribution plausible and varied? \textbf{R:} kinetic/geometric
and learned-metric subsets. Small GT cohorts limit FID interpretation. \\
& Velocity, acceleration, jerk, activity and repetition $\leftrightarrow$GT;
boundary jumps and physical violations $\downarrow$ &
\textbf{R:} motion-dynamics and continuity diagnostics. \textbf{I:} contact/skating
proxies and standardized physical checks. Low activity or low jerk alone is not success. \\
& Blinded naturalness, smoothness, expressiveness and aesthetics $\uparrow$ &
\textbf{P:} silent-video quality study and paired reference/execution study;
no perceptual superiority is claimed from existing proxy scores. \\
\bottomrule
\multicolumn{3}{@{}l@{}}{\footnotesize $\uparrow$: higher is better; $\downarrow$:
lower is better; $\leftrightarrow$GT: match the paired ground-truth distribution.} \\
\end{tabular}
\end{table*}

\subsection{Shared Evaluation Contract}
Table~\ref{tab:evaluation_matrix} separates three questions: learned cross-modal
matching (MMR), explicit music adaptation, and dance quality/aesthetics. They are
complementary, not independent quantities that can be added into one score.
The third block concerns the motion itself; the first two concern its relation
to the soundtrack. Distributional, rhythmic, physical, and perceptual evidence
follows established dance evaluation practice~\cite{li2021aist,tseng2023edge,li2024lodge,ji2026discoforcing,li2026roboperform}.

The calibration chain comprises paired human motion ($O_{\mathrm{Human}}$), its
GMR-retargeted G1 counterpart ($O_{\mathrm{G1}}$), and execution of that G1 target
($O_{\mathrm{Exec}}$). Generated motion is evaluated through the recorded
$M_{\mathrm{gen}}$--$M_{\mathrm{sent}}$--$M_{\mathrm{target}}$--$M_{\mathrm{exec}}$
chain. The generator is trained directly in the G1 domain; online output is not
retargeted again. The final protocol applies all three blocks to both
$M_{\mathrm{gen}}$ and $M_{\mathrm{exec}}$ where the evaluator's domain calibration
permits it. Existing tracking recordings do not imply that this complete scoring
matrix has already been computed.

Audio time, crop duration, rig, feature normalization, evaluator checkpoint, and
aggregation rules are fixed within a comparison. Music scores use the recorded
audio--action clock, without independently shifting each video to improve fit.
Lag-compensated tracking is a separate diagnostic. Startup is reported untrimmed
and separately from the post-startup interval. Missing inputs, unsupported
domains, and unvalidated endpoints are marked N/A, not treated as zero-quality
samples or silently removed. Music-conditioning features are not evaluation
metrics; the learned evaluator is frozen independently of generator fitting.

The source song, rather than the frame, is the statistical unit. Sampling seeds
are repeated measurements within a song, distinct from training seeds and
tracker resets. Final comparisons require a held-out manifest, all eligible
songs, paired song-level effects with 95\% confidence intervals, and tempo/style
strata based on verified metadata. Demonstration selection is separate from
quantitative inference. Different human/G1 feature spaces or evaluator versions
must not share a numerical ranking. Generated-to-GT joint RMSE is not an overall
choreography score because valid dance is one-to-many.

\subsection{Block 1: MMR Matching and Retrieval}
\textbf{Purpose and evaluator.} Music--motion retrieval (MMR) tests whether
corresponding audio and motion are closer than controlled mismatches in an
independently learned space. Our implemented evaluator uses 35-D Librosa
temporal audio descriptors and frozen MuQ-large global audio features with
native G1 motion features; its two towers
are fitted on paired training GT, then frozen for scoring. Global/local
contrastive and time-shift controls calibrate correspondence. Music identity,
not dancer or performance identity, defines positives. The model is an
evaluation instrument, not a new generator contribution or a substitute for
human judgment. Offline evaluation may inspect the complete scored segment;
these evaluator features are not supplied to the causal generator. This does
not grant future-waveform access to any deployable conditioning route.

\textbf{Paired matching.} Let $u_g,v_g$ be global music/motion embeddings and
$u_t,v_t$ their aligned, non-overlapping one-second segment embeddings. We use
the global--temporal matching distance adopted from SoulNet~\cite{li2025soulnet}:
\begin{align}
 D_{\mathrm{temp}} &= \sum_{t=2}^{T}
 \left\|(u_t-u_{t-1})-(v_t-v_{t-1})\right\|_2,\\
 d_{\mathrm{MMR}} &=
 \sqrt{0.7\|u_g-v_g\|_2^2+0.3D_{\mathrm{temp}}}.
 \label{eq:mmr_score}
\end{align}
MMR-MS averages this per-pair distance; lower is better. We also retain the
global and temporal terms to expose their different behavior. The temporal
term is a sum of norms, not squared norms or a time average; crop length must
therefore be matched. A global-only fallback is not the full MMR-MS.
MMDist separately measures mean paired global Euclidean distance; where
multiple matching candidates exist, the implementation uses their centroid.
It is not maximum mean discrepancy (MMD) and not a retrieval success rate.

\textbf{Retrieval.} We report audio-to-motion and motion-to-audio R@K for
$K\in\{1,2,3,5,10\}$, together with median and mean rank. A query succeeds if a
valid same-music candidate appears among its first $K$ results. Candidate pools,
query counts, distance convention, and multiple-positive rules remain fixed;
the first valid positive defines the rank. Chance comparisons must reflect
the actual pool and number of positives. Correct versus wrong-track and
wrong-time margins test whether the evaluator detects more than generic
movement, genre similarity, or performance identity.

\textbf{Status and limits.} FineDance-only matching/retrieval results are
available for formal generator subsets. A multi-dataset candidate has not
closed all AIST++ calibration gates, as detailed below. No cross-dataset or
executed-motion aesthetic ranking follows automatically from a lower MMR-MS.
Evaluator identity, training/test music separation, feature versions, and
temporal-embedding availability accompany each reported score.

\subsection{Block 2: Music Adaptation}
\textbf{Rhythmic alignment.} For motion-accent times $\mathcal{B}_x$ and
audio-beat times $\mathcal{B}_a$, motion-to-music Beat Alignment Score is
\begin{equation}
 \mathrm{BAS}=\frac{1}{|\mathcal{B}_x|}
 \sum_{b\in\mathcal{B}_x}
 \exp\!\left(-\frac{\min_{a\in\mathcal{B}_a}(b-a)^2}{2\sigma^2}\right).
 \label{eq:bas}
\end{equation}
The event extractor, smoothing, and $\sigma$ are frozen and recorded with the
metric version. The layer analyzer derives motion accents from minima of
smoothed mean joint speed and uses $\sigma=0.10$~s. Reverse-direction BAS is
labelled separately; historical scores require their own extractor/version
audit before pooling.
We retain event counts, beat density, nearest-event timing and offbeat
diagnostics: a high BAS with very few events does not establish sustained
rhythmic dance. Historical Beat-P/R/F1 values describe event coverage under a
specified matching tolerance; without an explicit beat-assignment protocol
they are not mandatory quality endpoints. Dance need not hit every audio beat,
and half-time or double-time phrasing is not inherently incorrect.

\textbf{Dynamics, tempo, and structure.} We measure zero-lag audio-onset
correlation with joint-speed and motion-impact curves. The implemented impact
proxy is the smoothed absolute derivative of the speed curve, not a measured
contact force. Best-lag correlation and its lag diagnose delayed response;
a lag from weak correlation is unreliable and is not used to retime the main
music score. Tempo/period and accent-phase diagnostics complement beat proximity
but require declared half/double-tempo handling. Phrase/section transition
response, annotated style fit, and affect agreement are further endpoints,
not inferred from BAS or motion energy. Validated phrase and affect evaluation
remains pending; human facial-expression metrics are not directly applicable
to G1's unactuated face.

\textbf{Appropriate adaptation rather than mere change.} Current formal
correct/null, wrong-track, wrong-time, and shifted-condition controls test
music dependence. The final cross-paired protocol additionally generates
$D_A,D_B$ from songs $A,B$ with matched motion history and sampling noise, and
scores both dances with both soundtracks. For a higher-is-better fit score $s$,
\begin{align}
 \Delta_{\mathrm{fit}}=\tfrac12\big[&s(A,D_A)-s(B,D_A)\nonumber\\
                              +&s(B,D_B)-s(A,D_B)\big].
 \label{eq:cross_music_fit}
\end{align}
For a distance metric, $s=-d$. Positive matched preference is sought together
with dance quality, not merely a large $\|D_A-D_B\|$. Same-style/similar-tempo
negatives test fine musical specificity; cross-style pairs test broader
response. This design is repeated for generated and executed dances without
changing their audio alignment. Blinded listeners rate rhythm, accent/dynamic
response, and overall music--dance fit; these perceptual controls remain TODO.

\subsection{Block 3: Dance Quality and Aesthetics}
\textbf{Distribution and diversity.} FID$_k$ and FID$_g$ compare generated and
held-out GT feature distributions using kinetic and geometric G1 descriptors,
respectively. FID$_L$ applies the same mean/covariance distance in an independent
frozen motion-embedding space. Div$_k$, Div$_g$, and Div$_L$ measure pairwise
feature distances and are interpreted relative to matched GT, not maximized
without bound. Same-song multimodality is the mean pairwise motion-embedding
distance across sampling seeds, macro-averaged over eligible songs; the formal
protocol uses at least three seeds. Here multimodality (MM) is distinct from
cross-modal MMR. Global diversity cannot show that one song admits diverse
valid dances. Small GT samples, including the current 18-song cohort, make FID
auxiliary evidence rather than a robust standalone superiority claim.

\textbf{Local and long-horizon motion.} We retain joint/root velocity,
acceleration and jerk distributions, high-frequency jitter, and position and
velocity discontinuities at C4 boundaries. Activity is measured using joint
excursion and squared angular velocity; static-window fraction and repeated
motion expose freezing or looping. These values are compared with GT and
controlled corruptions: standing still minimizes jerk but is not a good dance,
and high activity can be jitter. Reconstruction NRMSE, branch-swap sensitivity,
and continuous-code effective rank test the D+C mechanism, while the three
evaluation blocks test generated outcomes. They are not interchangeable.

\textbf{Physical plausibility.} Contact-weighted skating, penetration,
joint/dynamic-limit violations, root drift, and stability complement motion
statistics. The current G1 analyzer includes a physical-foot-contact (PFC)
proxy combining planar root acceleration and foot motion, and foot-skating
rate (FSR) proxies based on estimated contact. Their ground height, contact
threshold, speed threshold, and world/body frame must be disclosed; they are
not silently equated with standardized human-motion benchmarks. Kinematic
plausibility does not prove dynamic feasibility. Actual survival, falls, and
interventions are measured in controller execution.

\textbf{Perceptual quality.} A planned blinded silent-video study rates
naturalness, continuity/smoothness, expressiveness, and overall aesthetic
quality separately from the soundtrack-fit study. A second paired study asks
whether execution retains the reference's expressive movement. Presentation
matches rig, camera, duration, fps, and real playback speed, with randomized
order and participant/song-aware analysis. Low FID, low jerk, and high tracking
retention do not individually prove that dance is beautiful. No completed
human-quality result is claimed in this draft.

\subsection{Execution and Timing as Supporting Diagnostics}
Tracking and runtime measurements accompany, rather than replace, the three
quality blocks. We report unaligned and lag-aligned joint RMSE, root drift,
target-to-execution lag, joint-wise amplitude and velocity-power retention,
and limb/frequency-band attenuation. Startup and post-startup results are
separate. Audio cutoff, forecast publication/fetch, action deadlines, condition
age, valid FMS count, planning-time tails, deadline misses, holds/aborts, and
simulation real-time factor establish whether the intended online contract
was realized. Recording coverage and resource logs qualify comparisons;
successful retries alone cannot estimate reliability. A stable robot with
attenuated accents may track acceptably while scoring poorly on music fit.

\papertodo{T04/T08--T10: complete the three-block score matrix for generated and
executed motion; freeze event tolerances, windows, evaluator hashes, and
small-sample reporting. Run the cross-paired music-fit test and blinded aesthetic
study. Exploratory tempo/style/contact metrics and planned phrase/affect metrics
must not be presented as fully validated results.}



\section{Overlapping and Projector-Based Supervision}
\label{app:supervision}
The full-state comparison in Table~\ref{tab:decoupling_audit} changes a
supervision bundle, not the presence versus absence of training. The audited
legacy route applies weighted native reconstruction, derivatives, sparse FK,
contact, boundary, endpoint and residual penalties. Its q-only reconstruction
uses weighted Huber loss over all 34 native channels: root channels have weight
2.0, leg channels 1.5, waist channels 1.25, and the remaining arm channels 0.5.
Thus fine native articulation remains in the structural objective. That term
also includes native first-difference loss (0.5) and contact BCE (0.1), before
its outer coefficient of 0.5.

The frozen legacy trainer specifies full-reconstruction coefficients of 1.0
for native motion, 0.5 for velocity, 0.1 for acceleration, 0.5 for FK, 0.1 for
contact BCE, 0.2 for contact height, 0.1 for contact slide, 1.0 for VQ, and 0.2
for the amplitude guard. Boundary loss has outer weight 1.0 and internal
velocity/acceleration/jerk weights 1.0/0.5/0.25. Additional coefficients are
0.5 for the q-only endpoint-state loss, 0.25 for critical-state invariance,
and 0.01 for residual rate. Training corruption keeps the residual unchanged
with probability 0.55, removes it with probability 0.30, and adds channel-scaled
noise with probability 0.15. On this legacy path, corrupted reconstructions
can still be trained toward the full target.

In contrast, the proposed supervision uses only the four terms of
Eq.~\eqref{eq:codec_loss}. Full reconstruction uses normalized native motion,
boundary-aware velocity, all-link/sole geometry and geometry velocity.
The discrete and corrupted-residual reconstructions are supervised only on
the declared structure projector, excluding its exact detail complement.
Each projector loss equally averages its four normalized blocks. The outer
coefficients are 1.0 for full reconstruction, 1.0 for discrete structure,
0.25 for residual robustness, and 1.0 for VQ; the VQ commitment coefficient
is 0.25. The contact head and the separate legacy physical, endpoint,
invariance, amplitude, and residual-rate penalties are absent. Physical
diagnostics remain evaluation measures rather than hidden extra losses.

The source definitions and frozen experiment contract establish these
supervision differences. Exact checkpoint/launch-argument binding remains part
of the historical table provenance audit. This comparison supports the effect
of the supervision bundle; attributing it to any one removed penalty would
require an additional controlled ablation.


\bibliographystyle{IEEEtran}
\bibliography{Reference}
\end{document}
